\documentclass[tikz,border=10pt]{standalone}
\usepackage{ctex}
\usepackage{xcolor}
\usetikzlibrary{math}

\begin{document}
\begin{tikzpicture}[
  font=\sffamily,
]

\definecolor{红}{HTML}{E74C3C}
\definecolor{橙}{HTML}{F39C12}
\definecolor{黄}{HTML}{FFC116}
\definecolor{绿}{HTML}{52C41A}
\definecolor{蓝}{HTML}{3498DB}
\definecolor{紫}{HTML}{9D3DCF}
\definecolor{黑}{HTML}{0E1D69}
\definecolor{红L}{HTML}{FFE2D2}
\definecolor{橙L}{HTML}{FFFFA8}
\definecolor{黄L}{HTML}{FFFFAC}
\definecolor{绿L}{HTML}{E8FFB0}
\definecolor{蓝L}{HTML}{CAFFFF}
\definecolor{紫L}{HTML}{FFD3FF}
\definecolor{黑L}{HTML}{A4B3FF}
\definecolor{红LL}{HTML}{FFFFFF}
\definecolor{橙LL}{HTML}{FFFFDA}
\definecolor{黄LL}{HTML}{FFFFDE}
\definecolor{绿LL}{HTML}{FFFFE2}
\definecolor{蓝LL}{HTML}{FCFFFF}
\definecolor{紫LL}{HTML}{FFFFFF}
\definecolor{黑LL}{HTML}{D6E5FF}

  % === 中心圆：红色-入门-语法基础 ===
  \fill[红] (0,0) circle (1.5);
  \node[white, font=\bfseries\large, align=center] at (0,0) {\textbf{语法}\\基础};

  % === 橙色 - 普及- (r=1.5~3.0) ===
  % 搜索 - 无知识点
  \fill[橙L] (0.0:1.5) arc (0.0:51.42857142857143:1.5) -- (51.42857142857143:3.0) arc (51.42857142857143:0.0:3.0) -- cycle;
  \draw[橙, thin] (0.0:1.5) -- (0.0:3.0);
  % 动态规划 - 无知识点
  \fill[橙L] (51.42857142857143:1.5) arc (51.42857142857143:102.85714285714286:1.5) -- (102.85714285714286:3.0) arc (102.85714285714286:51.42857142857143:3.0) -- cycle;
  \draw[橙, thin] (51.42857142857143:1.5) -- (51.42857142857143:3.0);
  % 数据结构 - 无知识点
  \fill[橙L] (102.85714285714286:1.5) arc (102.85714285714286:154.28571428571428:1.5) -- (154.28571428571428:3.0) arc (154.28571428571428:102.85714285714286:3.0) -- cycle;
  \draw[橙, thin] (102.85714285714286:1.5) -- (102.85714285714286:3.0);
  % 图论 - 无知识点
  \fill[橙L] (154.28571428571428:1.5) arc (154.28571428571428:205.71428571428572:1.5) -- (205.71428571428572:3.0) arc (205.71428571428572:154.28571428571428:3.0) -- cycle;
  \draw[橙, thin] (154.28571428571428:1.5) -- (154.28571428571428:3.0);
  % 字符串 - 查找替换 大小写转换
  \fill[橙L] (205.71428571428572:1.5) arc (205.71428571428572:257.14285714285717:1.5) -- (257.14285714285717:3.0) arc (257.14285714285717:205.71428571428572:3.0) -- cycle;
  \node[\scriptsize, 橙!90!black, align=center, anchor=center, inner sep=1pt, text width=2.2cm] at (231.42857142857144:2.25) {\textbf{查找替换\\大小写转换}};
  \draw[橙, thin] (205.71428571428572:1.5) -- (205.71428571428572:3.0);
  % 数学 - 进制转换 高精度
  \fill[橙L] (257.14285714285717:1.5) arc (257.14285714285717:308.57142857142856:1.5) -- (308.57142857142856:3.0) arc (308.57142857142856:257.14285714285717:3.0) -- cycle;
  \node[\scriptsize, 橙!90!black, align=center, anchor=center, inner sep=1pt, text width=2.2cm] at (282.8571428571429:2.25) {\textbf{进制转换\\高精度}};
  \draw[橙, thin] (257.14285714285717:1.5) -- (257.14285714285717:3.0);
  % 计算几何 - 无知识点
  \fill[橙L] (308.57142857142856:1.5) arc (308.57142857142856:360.0:1.5) -- (360.0:3.0) arc (360.0:308.57142857142856:3.0) -- cycle;
  \draw[橙, thin] (308.57142857142856:1.5) -- (308.57142857142856:3.0);
  \draw[橙, thin] (0:3.0) arc (0:360:3.0);
  % === 黄色 - 普及/提高- (r=3.0~4.5) ===
  % 搜索 - DFS BFS 排列组合
  \fill[黄L] (0.0:3.0) arc (0.0:51.42857142857143:3.0) -- (51.42857142857143:4.5) arc (51.42857142857143:0.0:4.5) -- cycle;
  \node[\tiny, 黄!90!black, align=center, anchor=center, inner sep=1pt, text width=2.0cm] at (25.714285714285715:3.75) {\textbf{DFS\\BFS\\排列组合}};
  \draw[黄, thin] (0.0:3.0) -- (0.0:4.5);
  % 动态规划 - 线性DP 背包DP
  \fill[黄L] (51.42857142857143:3.0) arc (51.42857142857143:102.85714285714286:3.0) -- (102.85714285714286:4.5) arc (102.85714285714286:51.42857142857143:4.5) -- cycle;
  \node[\scriptsize, 黄!90!black, align=center, anchor=center, inner sep=1pt, text width=2.2cm] at (77.14285714285714:3.75) {\textbf{线性DP\\背包DP}};
  \draw[黄, thin] (51.42857142857143:3.0) -- (51.42857142857143:4.5);
  % 数据结构 - 栈 队列 并查集 set/map
  \fill[黄L] (102.85714285714286:3.0) arc (102.85714285714286:154.28571428571428:3.0) -- (154.28571428571428:4.5) arc (154.28571428571428:102.85714285714286:4.5) -- cycle;
  \node[\tiny, 黄!90!black, align=center, anchor=center, inner sep=1pt, text width=2.0cm] at (128.57142857142856:3.75) {\textbf{栈 队列\\并查集\\set/map}};
  \draw[黄, thin] (102.85714285714286:3.0) -- (102.85714285714286:4.5);
  % 图论 - 无知识点
  \fill[黄L] (154.28571428571428:3.0) arc (154.28571428571428:205.71428571428572:3.0) -- (205.71428571428572:4.5) arc (205.71428571428572:154.28571428571428:4.5) -- cycle;
  \draw[黄, thin] (154.28571428571428:3.0) -- (154.28571428571428:4.5);
  % 字符串 - 无知识点
  \fill[黄L] (205.71428571428572:3.0) arc (205.71428571428572:257.14285714285717:3.0) -- (257.14285714285717:4.5) arc (257.14285714285717:205.71428571428572:4.5) -- cycle;
  \draw[黄, thin] (205.71428571428572:3.0) -- (205.71428571428572:4.5);
  % 数学 - 素数 gcd/lcm 快速幂 组合数
  \fill[黄L] (257.14285714285717:3.0) arc (257.14285714285717:308.57142857142856:3.0) -- (308.57142857142856:4.5) arc (308.57142857142856:257.14285714285717:4.5) -- cycle;
  \node[\scriptsize, 黄!90!black, align=center, anchor=center, inner sep=1pt, text width=2.2cm] at (282.8571428571429:3.75) {\textbf{素数 gcd/lcm\\快速幂 组合数}};
  \draw[黄, thin] (257.14285714285717:3.0) -- (257.14285714285717:4.5);
  % 计算几何 - 无知识点
  \fill[黄L] (308.57142857142856:3.0) arc (308.57142857142856:360.0:3.0) -- (360.0:4.5) arc (360.0:308.57142857142856:4.5) -- cycle;
  \draw[黄, thin] (308.57142857142856:3.0) -- (308.57142857142856:4.5);
  \draw[黄, thin] (0:4.5) arc (0:360:4.5);
  % === 绿色 - 普及+/提高 (r=4.5~6.0) ===
  % 搜索 - 记忆化搜索 剪枝 双向BFS
  \fill[绿L] (0.0:4.5) arc (0.0:51.42857142857143:4.5) -- (51.42857142857143:6.0) arc (51.42857142857143:0.0:6.0) -- cycle;
  \node[\tiny, 绿!90!black, align=center, anchor=center, inner sep=1pt, text width=2.0cm] at (25.714285714285715:5.25) {\textbf{记忆化搜索\\剪枝\\双向BFS}};
  \draw[绿, thin] (0.0:4.5) -- (0.0:6.0);
  % 动态规划 - 区间DP 树形DP 状压入门 数位DP
  \fill[绿L] (51.42857142857143:4.5) arc (51.42857142857143:102.85714285714286:4.5) -- (102.85714285714286:6.0) arc (102.85714285714286:51.42857142857143:6.0) -- cycle;
  \node[\tiny, 绿!90!black, align=center, anchor=center, inner sep=1pt, text width=2.0cm] at (77.14285714285714:5.25) {\textbf{区间DP 树形DP\\状压入门\\数位DP}};
  \draw[绿, thin] (51.42857142857143:4.5) -- (51.42857142857143:6.0);
  % 数据结构 - 线段树(单点) 树状数组 ST表 分块
  \fill[绿L] (102.85714285714286:4.5) arc (102.85714285714286:154.28571428571428:4.5) -- (154.28571428571428:6.0) arc (154.28571428571428:102.85714285714286:6.0) -- cycle;
  \node[\tiny, 绿!90!black, align=center, anchor=center, inner sep=1pt, text width=2.0cm] at (128.57142857142856:5.25) {\textbf{线段树(单点)\\树状数组\\ST表 分块}};
  \draw[绿, thin] (102.85714285714286:4.5) -- (102.85714285714286:6.0);
  % 图论 - Dijkstra Floyd Kruskal Prim 拓扑 二分图
  \fill[绿L] (154.28571428571428:4.5) arc (154.28571428571428:205.71428571428572:4.5) -- (205.71428571428572:6.0) arc (205.71428571428572:154.28571428571428:6.0) -- cycle;
  \node[\tiny, 绿!90!black, align=center, anchor=center, inner sep=1pt, text width=2.0cm] at (180.0:5.25) {\textbf{Dijkstra Floyd\\Kruskal Prim\\拓扑 二分图}};
  \draw[绿, thin] (154.28571428571428:4.5) -- (154.28571428571428:6.0);
  % 字符串 - KMP Trie 字符串哈希
  \fill[绿L] (205.71428571428572:4.5) arc (205.71428571428572:257.14285714285717:4.5) -- (257.14285714285717:6.0) arc (257.14285714285717:205.71428571428572:6.0) -- cycle;
  \node[\tiny, 绿!90!black, align=center, anchor=center, inner sep=1pt, text width=2.0cm] at (231.42857142857144:5.25) {\textbf{KMP\\Trie\\字符串哈希}};
  \draw[绿, thin] (205.71428571428572:4.5) -- (205.71428571428572:6.0);
  % 数学 - exgcd CRT Lucas 高斯消元 欧拉函数
  \fill[绿L] (257.14285714285717:4.5) arc (257.14285714285717:308.57142857142856:4.5) -- (308.57142857142856:6.0) arc (308.57142857142856:257.14285714285717:6.0) -- cycle;
  \node[\tiny, 绿!90!black, align=center, anchor=center, inner sep=1pt, text width=2.0cm] at (282.8571428571429:5.25) {\textbf{exgcd CRT\\Lucas 高斯消元\\欧拉函数}};
  \draw[绿, thin] (257.14285714285717:4.5) -- (257.14285714285717:6.0);
  % 计算几何 - 凸包(Graham) 叉积/点积
  \fill[绿L] (308.57142857142856:4.5) arc (308.57142857142856:360.0:4.5) -- (360.0:6.0) arc (360.0:308.57142857142856:6.0) -- cycle;
  \node[\scriptsize, 绿!90!black, align=center, anchor=center, inner sep=1pt, text width=2.2cm] at (334.2857142857143:5.25) {\textbf{凸包(Graham)\\叉积/点积}};
  \draw[绿, thin] (308.57142857142856:4.5) -- (308.57142857142856:6.0);
  \draw[绿, thin] (0:6.0) arc (0:360:6.0);
  % === 蓝色 - 提高+/省选- (r=6.0~7.5) ===
  % 搜索 - 无知识点
  \fill[蓝L] (0.0:6.0) arc (0.0:51.42857142857143:6.0) -- (51.42857142857143:7.5) arc (51.42857142857143:0.0:7.5) -- cycle;
  \draw[蓝, thin] (0.0:6.0) -- (0.0:7.5);
  % 动态规划 - 状压DP 单调队列优化 斜率优化入门
  \fill[蓝L] (51.42857142857143:6.0) arc (51.42857142857143:102.85714285714286:6.0) -- (102.85714285714286:7.5) arc (102.85714285714286:51.42857142857143:7.5) -- cycle;
  \node[\tiny, 蓝!90!black, align=center, anchor=center, inner sep=1pt, text width=2.0cm] at (77.14285714285714:6.75) {\textbf{状压DP\\单调队列优化\\斜率优化入门}};
  \draw[蓝, thin] (51.42857142857143:6.0) -- (51.42857142857143:7.5);
  % 数据结构 - 平衡树 树剖 主席树(可持久化) 莫队 CDQ分治
  \fill[蓝L] (102.85714285714286:6.0) arc (102.85714285714286:154.28571428571428:6.0) -- (154.28571428571428:7.5) arc (154.28571428571428:102.85714285714286:7.5) -- cycle;
  \node[\tiny, 蓝!90!black, align=center, anchor=center, inner sep=1pt, text width=2.0cm] at (128.57142857142856:6.75) {\textbf{平衡树 树剖\\主席树(可持久化)\\莫队 CDQ分治}};
  \draw[蓝, thin] (102.85714285714286:6.0) -- (102.85714285714286:7.5);
  % 图论 - 网络流(Dinic) 费用流 Tarjan 2-SAT 差分约束
  \fill[蓝L] (154.28571428571428:6.0) arc (154.28571428571428:205.71428571428572:6.0) -- (205.71428571428572:7.5) arc (205.71428571428572:154.28571428571428:7.5) -- cycle;
  \node[\tiny, 蓝!90!black, align=center, anchor=center, inner sep=1pt, text width=2.0cm] at (180.0:6.75) {\textbf{网络流(Dinic)\\费用流 Tarjan\\2-SAT 差分约束}};
  \draw[蓝, thin] (154.28571428571428:6.0) -- (154.28571428571428:7.5);
  % 字符串 - AC自动机 Manacher 后缀数组 SA SAM入门
  \fill[蓝L] (205.71428571428572:6.0) arc (205.71428571428572:257.14285714285717:6.0) -- (257.14285714285717:7.5) arc (257.14285714285717:205.71428571428572:7.5) -- cycle;
  \node[\tiny, 蓝!90!black, align=center, anchor=center, inner sep=1pt, text width=2.0cm] at (231.42857142857144:6.75) {\textbf{AC自动机\\Manacher\\后缀数组 SA\\SAM入门}};
  \draw[蓝, thin] (205.71428571428572:6.0) -- (205.71428571428572:7.5);
  % 数学 - FFT/NTT 线性基 莫比乌斯反演 杜教筛 SG函数
  \fill[蓝L] (257.14285714285717:6.0) arc (257.14285714285717:308.57142857142856:6.0) -- (308.57142857142856:7.5) arc (308.57142857142856:257.14285714285717:7.5) -- cycle;
  \node[\tiny, 蓝!90!black, align=center, anchor=center, inner sep=1pt, text width=2.0cm] at (282.8571428571429:6.75) {\textbf{FFT/NTT\\线性基\\莫比乌斯反演\\杜教筛 SG函数}};
  \draw[蓝, thin] (257.14285714285717:6.0) -- (257.14285714285717:7.5);
  % 计算几何 - 半平面交 旋转卡壳
  \fill[蓝L] (308.57142857142856:6.0) arc (308.57142857142856:360.0:6.0) -- (360.0:7.5) arc (360.0:308.57142857142856:7.5) -- cycle;
  \node[\scriptsize, 蓝!90!black, align=center, anchor=center, inner sep=1pt, text width=2.2cm] at (334.2857142857143:6.75) {\textbf{半平面交\\旋转卡壳}};
  \draw[蓝, thin] (308.57142857142856:6.0) -- (308.57142857142856:7.5);
  \draw[蓝, thin] (0:7.5) arc (0:360:7.5);
  % === 紫色 - 省选/NOI- (r=7.5~9.0) ===
  % 搜索 - 无知识点
  \fill[紫L] (0.0:7.5) arc (0.0:51.42857142857143:7.5) -- (51.42857142857143:9.0) arc (51.42857142857143:0.0:9.0) -- cycle;
  \draw[紫, thin] (0.0:7.5) -- (0.0:9.0);
  % 动态规划 - 插头DP 动态DP 决策单调性 斯坦纳树
  \fill[紫L] (51.42857142857143:7.5) arc (51.42857142857143:102.85714285714286:7.5) -- (102.85714285714286:9.0) arc (102.85714285714286:51.42857142857143:9.0) -- cycle;
  \node[\tiny, 紫!90!black, align=center, anchor=center, inner sep=1pt, text width=2.0cm] at (77.14285714285714:8.25) {\textbf{插头DP\\动态DP\\决策单调性\\斯坦纳树}};
  \draw[紫, thin] (51.42857142857143:7.5) -- (51.42857142857143:9.0);
  % 数据结构 - LCT(动态树) 树套树 KD-Tree 点分治 虚树
  \fill[紫L] (102.85714285714286:7.5) arc (102.85714285714286:154.28571428571428:7.5) -- (154.28571428571428:9.0) arc (154.28571428571428:102.85714285714286:9.0) -- cycle;
  \node[\tiny, 紫!90!black, align=center, anchor=center, inner sep=1pt, text width=2.0cm] at (128.57142857142856:8.25) {\textbf{LCT(动态树)\\树套树 KD-Tree\\点分治 虚树}};
  \draw[紫, thin] (102.85714285714286:7.5) -- (102.85714285714286:9.0);
  % 图论 - 上下界网络流 最大权闭合子图 最小树形图 支配树
  \fill[紫L] (154.28571428571428:7.5) arc (154.28571428571428:205.71428571428572:7.5) -- (205.71428571428572:9.0) arc (205.71428571428572:154.28571428571428:9.0) -- cycle;
  \node[\tiny, 紫!90!black, align=center, anchor=center, inner sep=1pt, text width=2.0cm] at (180.0:8.25) {\textbf{上下界网络流\\最大权闭合子图\\最小树形图 支配树}};
  \draw[紫, thin] (154.28571428571428:7.5) -- (154.28571428571428:9.0);
  % 字符串 - SAM进阶 回文自动机 PAM 广义SAM
  \fill[紫L] (205.71428571428572:7.5) arc (205.71428571428572:257.14285714285717:7.5) -- (257.14285714285717:9.0) arc (257.14285714285717:205.71428571428572:9.0) -- cycle;
  \node[\tiny, 紫!90!black, align=center, anchor=center, inner sep=1pt, text width=2.0cm] at (231.42857142857144:8.25) {\textbf{SAM进阶\\回文自动机 PAM\\广义SAM}};
  \draw[紫, thin] (205.71428571428572:7.5) -- (205.71428571428572:9.0);
  % 数学 - 多项式全家桶 FWT Burnside/Polya 二次剩余
  \fill[紫L] (257.14285714285717:7.5) arc (257.14285714285717:308.57142857142856:7.5) -- (308.57142857142856:9.0) arc (308.57142857142856:257.14285714285717:9.0) -- cycle;
  \node[\tiny, 紫!90!black, align=center, anchor=center, inner sep=1pt, text width=2.0cm] at (282.8571428571429:8.25) {\textbf{多项式全家桶\\FWT\\Burnside/Polya\\二次剩余}};
  \draw[紫, thin] (257.14285714285717:7.5) -- (257.14285714285717:9.0);
  % 计算几何 - 圆的面积并 三维几何 动态凸包
  \fill[紫L] (308.57142857142856:7.5) arc (308.57142857142856:360.0:7.5) -- (360.0:9.0) arc (360.0:308.57142857142856:9.0) -- cycle;
  \node[\tiny, 紫!90!black, align=center, anchor=center, inner sep=1pt, text width=2.0cm] at (334.2857142857143:8.25) {\textbf{圆的面积并\\三维几何\\动态凸包}};
  \draw[紫, thin] (308.57142857142856:7.5) -- (308.57142857142856:9.0);
  \draw[紫, thin] (0:9.0) arc (0:360:9.0);
  % === 黑色 - NOI/NOI+/CTSC (r=9.0~10.5) ===
  % 搜索 - Meet-in-the- middle
  \fill[黑L] (0.0:9.0) arc (0.0:51.42857142857143:9.0) -- (51.42857142857143:10.5) arc (51.42857142857143:0.0:10.5) -- cycle;
  \node[\scriptsize, 黑!90!black, align=center, anchor=center, inner sep=1pt, text width=2.2cm] at (25.714285714285715:9.75) {\textbf{Meet-in-the-\\middle}};
  \draw[黑, thin] (0.0:9.0) -- (0.0:10.5);
  % 动态规划 - 论文级DP
  \fill[黑L] (51.42857142857143:9.0) arc (51.42857142857143:102.85714285714286:9.0) -- (102.85714285714286:10.5) arc (102.85714285714286:51.42857142857143:10.5) -- cycle;
  \node[\footnotesize, 黑!90!black, align=center, anchor=center, inner sep=1pt, text width=2.2cm] at (77.14285714285714:9.75) {\textbf{论文级DP}};
  \draw[黑, thin] (51.42857142857143:9.0) -- (51.42857142857143:10.5);
  % 数据结构 - ETT 树分块
  \fill[黑L] (102.85714285714286:9.0) arc (102.85714285714286:154.28571428571428:9.0) -- (154.28571428571428:10.5) arc (154.28571428571428:102.85714285714286:10.5) -- cycle;
  \node[\scriptsize, 黑!90!black, align=center, anchor=center, inner sep=1pt, text width=2.2cm] at (128.57142857142856:9.75) {\textbf{ETT\\树分块}};
  \draw[黑, thin] (102.85714285714286:9.0) -- (102.85714285714286:10.5);
  % 图论 - 一般图匹配 (带花树) 平面图算法
  \fill[黑L] (154.28571428571428:9.0) arc (154.28571428571428:205.71428571428572:9.0) -- (205.71428571428572:10.5) arc (205.71428571428572:154.28571428571428:10.5) -- cycle;
  \node[\tiny, 黑!90!black, align=center, anchor=center, inner sep=1pt, text width=2.0cm] at (180.0:9.75) {\textbf{一般图匹配\\(带花树)\\平面图算法}};
  \draw[黑, thin] (154.28571428571428:9.0) -- (154.28571428571428:10.5);
  % 字符串 - 字符串+ 生成函数
  \fill[黑L] (205.71428571428572:9.0) arc (205.71428571428572:257.14285714285717:9.0) -- (257.14285714285717:10.5) arc (257.14285714285717:205.71428571428572:10.5) -- cycle;
  \node[\scriptsize, 黑!90!black, align=center, anchor=center, inner sep=1pt, text width=2.2cm] at (231.42857142857144:9.75) {\textbf{字符串+\\生成函数}};
  \draw[黑, thin] (205.71428571428572:9.0) -- (205.71428571428572:10.5);
  % 数学 - 多项式多点求值 快速插值 BM算法
  \fill[黑L] (257.14285714285717:9.0) arc (257.14285714285717:308.57142857142856:9.0) -- (308.57142857142856:10.5) arc (308.57142857142856:257.14285714285717:10.5) -- cycle;
  \node[\tiny, 黑!90!black, align=center, anchor=center, inner sep=1pt, text width=2.0cm] at (282.8571428571429:9.75) {\textbf{多项式多点求值\\快速插值\\BM算法}};
  \draw[黑, thin] (257.14285714285717:9.0) -- (257.14285714285717:10.5);
  % 计算几何 - 三维凸包
  \fill[黑L] (308.57142857142856:9.0) arc (308.57142857142856:360.0:9.0) -- (360.0:10.5) arc (360.0:308.57142857142856:10.5) -- cycle;
  \node[\footnotesize, 黑!90!black, align=center, anchor=center, inner sep=1pt, text width=2.2cm] at (334.2857142857143:9.75) {\textbf{三维凸包}};
  \draw[黑, thin] (308.57142857142856:9.0) -- (308.57142857142856:10.5);
  \draw[黑, thin] (0:10.5) arc (0:360:10.5);
  \draw[black!30, thin] (0.0:1.5) -- (0.0:10.5);
  \draw[black!30, thin] (51.42857142857143:1.5) -- (51.42857142857143:10.5);
  \draw[black!30, thin] (102.85714285714286:1.5) -- (102.85714285714286:10.5);
  \draw[black!30, thin] (154.28571428571428:1.5) -- (154.28571428571428:10.5);
  \draw[black!30, thin] (205.71428571428572:1.5) -- (205.71428571428572:10.5);
  \draw[black!30, thin] (257.14285714285717:1.5) -- (257.14285714285717:10.5);
  \draw[black!30, thin] (308.57142857142856:1.5) -- (308.57142857142856:10.5);

  % === 扇区标签 ===
  \node[font=\small\bfseries, align=center] at (25.714285714285715:11.1) {搜索};
  \node[font=\small\bfseries, align=center] at (77.14285714285714:11.1) {动态规划};
  \node[font=\small\bfseries, align=center] at (128.57142857142858:11.1) {数据结构};
  \node[font=\small\bfseries, align=center] at (180.0:11.1) {图论};
  \node[font=\small\bfseries, align=center] at (231.42857142857142:11.1) {字符串};
  \node[font=\small\bfseries, align=center] at (282.85714285714283:11.1) {数学};
  \node[font=\small\bfseries, align=center] at (334.2857142857143:11.1) {计算几何};

  % === 难度标签(右侧) ===
  \node[黑, font=\scriptsize, anchor=west] at (0:9.85) {\textcolor{黑}{\rule{4pt}{4pt}} NOI/NOI+/CTSC};
  \node[紫, font=\scriptsize, anchor=west] at (0:8.35) {\textcolor{紫}{\rule{4pt}{4pt}} 省选/NOI-};
  \node[蓝, font=\scriptsize, anchor=west] at (0:6.85) {\textcolor{蓝}{\rule{4pt}{4pt}} 提高+/省选-};
  \node[绿, font=\scriptsize, anchor=west] at (0:5.35) {\textcolor{绿}{\rule{4pt}{4pt}} 普及+/提高};
  \node[黄, font=\scriptsize, anchor=west] at (0:3.85) {\textcolor{黄}{\rule{4pt}{4pt}} 普及/提高-};
  \node[橙, font=\scriptsize, anchor=west] at (0:2.35) {\textcolor{橙}{\rule{4pt}{4pt}} 普及-};
  \node[红, font=\scriptsize, anchor=west] at (0:0.85) {\textcolor{红}{\rule{4pt}{4pt}} 入门};
\end{tikzpicture}
\end{document}
